\documentclass[11pt]{amsart}
\usepackage{amsmath,amssymb,amsthm}
\usepackage[margin=1.2in]{geometry}
\usepackage{newtxtext,newtxmath}
\usepackage{booktabs}
\usepackage{microtype}
\usepackage[colorlinks=true, linkcolor=blue, citecolor=blue, urlcolor=blue]{hyperref}

\newtheorem{theorem}{Theorem}[section]
\newtheorem{lemma}[theorem]{Lemma}
\newtheorem{corollary}[theorem]{Corollary}
\newtheorem{proposition}[theorem]{Proposition}
\newtheorem{remark}[theorem]{Remark}
\newtheorem{hypothesis}[theorem]{Hypothesis}

\newcommand{\Z}{\mathbb{Z}}
\newcommand{\llog}{\log\log}
\newcommand{\lllog}{\log\log\log}
\newcommand{\CC}{C}

\begin{document}

\title[Primes near Carmichael numbers]{Primes near Carmichael numbers:\\ a trade-off theorem and its barriers}

\author{Dan Brown$^*$}
\email{jdb1729@gmail.com}
\thanks{$^*$Author of \emph{The Da Vinci Code}.}

\date{August 24, 2026}

\begin{abstract}
Let $\CC(X)$ denote the number of Carmichael numbers up to $X$. We prove an
unconditional trade-off: if $\CC(X) - \CC(X/2) \ge X^{\theta}$, then some
Carmichael number $m \in (X/2, X]$ has a prime within $m^{1.23 - \theta +
\varepsilon}$ of it. Consequently any counting bound $\CC(X) \ge X^{0.706}$
would yield primes nearer to (some) Carmichael numbers than the
Baker--Harman--Pintz exponent $0.525$ guarantees near arbitrary points, with no
new prime-gap input; under the Erd\H{o}s--Pomerance abundance conjecture the
exponent drops to $0.23 + \varepsilon$, and under the Lindel\"of hypothesis to
$\varepsilon$. We complement this with two negative results. First, a barrier:
since $\CC(X) \le X^{1 - (1+o(1))\lllog X/\llog X}$ unconditionally, no
pigeonhole of this type --- under any abundance assumption and any
variance-strength bound on primeless neighborhoods, including those available
under the Riemann Hypothesis --- can place a prime within
$\exp\bigl((1-\varepsilon)\,\lllog m \cdot \log m/\llog m\bigr)$ of a Carmichael
number; in particular polylogarithmic proximity is out of reach for the method. Second, an
algorithmic accounting: a deterministic prime-finding algorithm whose
correctness certificate is a counting inequality of this type must test
$X^{1+o(1)}$ candidates, so it cannot compete with the trivial
$X^{0.525+o(1)}$ scan, let alone the $X^{1/2+o(1)}$ record of Lagarias and
Odlyzko. On the positive side, the Alford--Granville--Pomerance
equidistribution theorem underlying the companion construction is itself a
prime-existence theorem for a thin explicit set: it yields, unconditionally
and effectively, a deterministic algorithm finding a prime in $(N/2, N]$ in
time $N^{7/12+o(1)}$ --- to our knowledge new, and the
third-fastest proven method. Its exponent is $1 - 1/A$ for any admissible
zero-density exponent $A$, hence $N^{1/2+\varepsilon}$ under the density
hypothesis and never below $N^{1/2}$, since $A \le 2$ is impossible.
Empirically, Carmichael numbers are not measurably
closer to primes than random odd numbers of the same size, although the
primes near them are confined to a sparse, predictable set of admissible
offsets.
\end{abstract}

\maketitle

\section{Introduction}
\label{sec:intro}

A companion paper \cite{Bro} constructs, deterministically and unconditionally,
a Carmichael number exceeding a given bound $n$ in time quasi-polylogarithmic in
$n$. By contrast, the fastest deterministic algorithm known to produce a
\emph{prime} exceeding $n$ runs in time $n^{1/2+o(1)}$ \cite{LO, Pol}. This
prompts a natural question: \emph{must primes lie close to
Carmichael numbers?} If Carmichael numbers --- which we can find quickly ---
were provably accompanied by nearby primes, one could find primes by searching
their neighborhoods, and any proximity exponent below $1/2$ would beat the
record.

This note is an after-action report on the most natural proof strategy, a
pigeonhole against prime-gap moments: what it delivers, and where it is
stopped. Write $p_1 < p_2 < \cdots$ for the
primes, $d_n = p_{n+1} - p_n$, and let
\[
E_H(X) \;=\; \{\, x \in (X/2,\, X] : \text{no prime in } [x - H,\, x + H] \,\}.
\]
Write $\CC(X)$ for the number of Carmichael numbers up to $X$, and
$\ell_2 = \llog X$, $\ell_3 = \lllog X$.

Our inputs are classical and recent bounds on both sides of the pigeonhole.
For the gaps: Selberg \cite{Sel} proved, under the Riemann Hypothesis, that
$\sum_{p_n \le x} d_n^2 \ll x \log^3 x$; unconditionally, Heath-Brown
\cite{HB} obtained the exponent $23/18$, improved through work of Peck and
Maynard to $5/4$, and by Stadlmann \cite{Sta} to
\begin{equation}
\label{eq:sta}
\sum_{p_n \le x} d_n^2 \;\ll_\varepsilon\; x^{1.23 + \varepsilon};
\end{equation}
under the Lindel\"of hypothesis Yu obtained exponent $1 + \varepsilon$ (see the
introduction of \cite{Sta} for this history). For the Carmichael counts:
$\CC(X) \ge X^{1/3}$ for large $X$ (Harman \cite{Har}, refining
Alford--Granville--Pomerance \cite{AGP}), while unconditionally
\begin{equation}
\label{eq:pom}
\CC(X) \;\le\; X^{\,1 - (1+o(1))\,\ell_3/\ell_2}
\end{equation}
(Pomerance \cite{Pom}), and the standard conjecture of Erd\H{o}s and Pomerance
(see \cite{Erd, GP}) asserts that \eqref{eq:pom} is an equality.

\section{The trade-off theorem}
\label{sec:tradeoff}

\begin{lemma}
\label{lem:E}
For $1 \le H \le X/8$ and every $\varepsilon > 0$,
\[
|E_H(X)| \;\le\; \frac{1}{2H} \sum_{p_n \le 2X} d_n^2
\;\ll_\varepsilon\; \frac{X^{1.23+\varepsilon}}{H}.
\]
Under the Riemann Hypothesis, if additionally $H \ge 12 \log X$, then
$|E_H(X)| \ll X \log^3 X / H$.
\end{lemma}

\begin{proof}
If $x \in E_H(X)$ then the consecutive primes $p_n < x < p_{n+1}$ that bracket
$x$ satisfy $p_{n+1} - x > H$ and $x - p_n > H$, so $d_n > 2H$; moreover
$p_n < X$ and $p_{n+1} \le x + d_n \le 2X$ once $d_n \le X$ (gaps below $2X$
of length $> X$ do not exist for large $X$ by Bertrand's postulate). Each such
gap contains fewer than $d_n$ integers $x$, whence
$|E_H(X)| \le \sum_{d_n > 2H,\, p_n \le 2X} d_n \le (2H)^{-1} \sum_{p_n \le 2X}
d_n^2$, and \eqref{eq:sta} gives the first bound.

For the second: Selberg's theorem \cite{Sel} gives, under RH,
$\int_{X/2}^{X} \bigl(\psi(x+H) - \psi(x) - H\bigr)^2\, dx \ll H X \log^2 X$
(we use the form with a $\log^3$-safe constant; any bound
$\ll HX\log^{O(1)}X$ serves). If $x \in E_H(X)$ then the interval
$(x, x+H]$ contains no primes, only prime powers $p^j$, $j \ge 2$, of which
there are at most $H X^{-1/2} + 2$, each contributing at most $\log 2X$ to
$\psi$; thus $\psi(x+H) - \psi(x) \le (H X^{-1/2} + 2)\log 2X \le H/2$ for
$H \ge 12 \log X$ and $X$ large, so the integrand is at least $H^2/4$ on
$E_H(X)$. Comparing measures gives $|E_H(X)| \ll X \log^3 X / H$.
\end{proof}

\begin{theorem}
\label{thm:tradeoff}
Let $0 < \theta \le 1$ and $\varepsilon > 0$, and suppose
$\CC(X) - \CC(X/2) \ge X^{\theta}$ for all large $X$. Then there are
infinitely many Carmichael numbers $m$ admitting a prime $p$ with
\[
|p - m| \;\le\; m^{\,1.23 - \theta + \varepsilon}.
\]
\end{theorem}

\begin{proof}
Set $H = X^{1.23 - \theta + \varepsilon}$ (if the exponent is $\ge 1$ the claim
is trivial by Bertrand's postulate, so assume it is $< 1$). By
Lemma~\ref{lem:E}, $|E_H(X)| \ll X^{1.23 + \varepsilon/2}/H =
X^{\theta - \varepsilon/2} < X^{\theta} \le \CC(X) - \CC(X/2)$ for large $X$.
Hence some Carmichael $m \in (X/2, X]$ lies outside $E_H(X)$ and has a prime
within $H = X^{1.23-\theta+\varepsilon} \le (2m)^{1.23-\theta+\varepsilon} \le
m^{1.23-\theta+2\varepsilon}$ for large $X$; taking $X = 2^j \to \infty$
produces infinitely many such $m$, and relabelling $\varepsilon$ finishes.
\end{proof}

\begin{corollary}[unconditional threshold]
\label{cor:threshold}
If $\CC(X) - \CC(X/2) \ge X^{\theta}$ with $\theta > 1.23 - 0.525 = 0.705$,
then infinitely many Carmichael numbers have primes within
$m^{1.23 - \theta + \varepsilon}$, closer than the exponent $0.525$
that Baker--Harman--Pintz \cite{BHP} guarantees for arbitrary points.
\end{corollary}

The current record is $\theta = 1/3$ \cite{Har}, so
Corollary~\ref{cor:threshold} is presently empty: for $\theta \le 0.705$ the
theorem returns exponents worse than $0.525$. Its content is a
\emph{reduction}: closing the gap from $0.33$ to $0.706$ in Carmichael
counting --- a problem entirely on the composite side, where the
Alford--Granville--Pomerance machinery keeps improving --- would
automatically produce primes abnormally close to Carmichael numbers, with
prime-gap technology frozen at \eqref{eq:sta}.

\begin{corollary}[under abundance]
\label{cor:abundance}
Assume $\CC(X) - \CC(X/2) \ge X^{1-\delta}$ for every $\delta > 0$ (a weak
form of the Erd\H{o}s--Pomerance conjecture). Then for every
$\varepsilon > 0$ there are infinitely many Carmichael $m$ with a prime
within $m^{0.23 + \varepsilon}$. Under the Lindel\"of hypothesis the exponent
improves to $\varepsilon$, for every $\varepsilon > 0$.
\end{corollary}

\begin{corollary}[under RH and abundance]
\label{cor:rh}
Assume RH and the asymptotic form
$\CC(X) - \CC(X/2) \ge X \exp\bigl(-(1+\varepsilon)\,\ell_3 \log X/\ell_2\bigr)$
of the Erd\H{o}s--Pomerance conjecture. Then infinitely many Carmichael $m$
have a prime within
$\exp\bigl((1+2\varepsilon)\,\lllog m \cdot \log m/\llog m\bigr) = m^{o(1)}$.
\end{corollary}

\begin{proof}
Take $H = \log^3 X \cdot \exp\bigl((1+\varepsilon)\ell_3 \log X/\ell_2\bigr)$
in the RH part of Lemma~\ref{lem:E}.
\end{proof}

\section{Two barriers}
\label{sec:barriers}

The exponents of Corollaries~\ref{cor:abundance}--\ref{cor:rh} are the floor
of the method.

\begin{proposition}[pigeonhole barrier]
\label{prop:barrier}
Let $B(X, H)$ be any provable upper bound for $|E_H(X)|$ of the shape
$B(X,H) \ge X (\log X)^{-A}$ for fixed $A$ whenever
$H \le (\log X)^{A}$ --- as is the case for every bound derived from a
variance (second-moment) estimate for primes in short intervals, including
Selberg's RH bound. Then the inequality $B(X,H) < \CC(X)$ required by the
pigeonhole fails for all such $H$, unconditionally: by \eqref{eq:pom},
$\CC(X) \le X^{1-(1+o(1))\ell_3/\ell_2} < X(\log X)^{-A}$ for every fixed $A$
and large $X$. Consequently no pigeonhole against Carmichael counts --- with
\emph{any} strength of abundance input, since \eqref{eq:pom} is an upper
bound --- can certify a prime within polylogarithmic distance of a Carmichael
number; the least $H$ it can certify is
$\exp\bigl((1-o(1))\,\ell_3 \log X/\ell_2\bigr)$, matching
Corollary~\ref{cor:rh}.
\end{proposition}

\begin{proof}
Immediate from \eqref{eq:pom} and the displayed shapes: for
$H \le (\log X)^A$ one has $B(X,H) \ge X(\log X)^{-A} > \CC(X)$, so the
pigeonhole cannot close. The floor statement follows by solving
$X \log^{O(1)} X/H < X^{1-(1+o(1))\ell_3/\ell_2}$ for $H$.
\end{proof}

Thus the interesting regime of Theorem~\ref{thm:tradeoff} is polynomial
proximity $m^{c}$; polylogarithmic proximity is denied territory for
counting methods, and entry requires either gap-tail
bounds far beyond second moments (Cram\'er-model strength: the obstruction
identified by the Polymath project \cite{Pol}), or
structural information about Carmichael positions, such as equidistribution of
subset products of the \cite{AGP} prime pool in short intervals, which lies
beyond current exponential-sum technology.

The second barrier concerns the intended algorithmic use.

\begin{proposition}[cost of counting certificates]
\label{prop:cost}
Consider a deterministic algorithm which outputs a prime by testing candidates
drawn from $R$ neighborhoods of radius $H$ centred at points of $(X/2, X]$
(Carmichael or otherwise), and whose correctness proof is the counting
inequality $R > B(X,H)$ with $B$ a provable bound for $|E_H(X)|$ as above.
In the worst case the algorithm performs at least $B(X,H) \cdot H / \log X$
primality tests. With the unconditional bound \eqref{eq:sta} this is
$\gg X^{1.23 - o(1)}$; with Selberg's RH bound it is $\gg X (\log X)^{2}$.
In both cases the cost exceeds $X$, so such an algorithm is slower than the
trivial scan of a single interval of length $X^{0.525}$ \cite{BHP}, and far
slower than the $X^{1/2+o(1)}$ algorithm of Lagarias--Odlyzko \cite{LO}.
\end{proposition}

\begin{proof}
The certificate guarantees only that the \emph{union} of the $R$
neighborhoods meets the primes; it identifies no neighborhood in particular,
and any $R - 1$ of them may be barren, since $B$ is the best available bound
on how many barren centres exist and $R$ must exceed it. An adversarial
configuration consistent with the certificate places the guaranteed prime
among the last candidates examined, whichever order the algorithm chooses, so
the worst case examines a constant fraction of the
$\asymp R H/\log X$ prime-candidates in the union: at least
$B(X,H)\, H/\log X$ tests. The two instantiations are
$B \cdot H \gg X^{1.23-o(1)}$ and $B \cdot H \gg X \log^{3} X$.
\end{proof}

Proposition~\ref{prop:cost} is a strong negative for the motivating program:
\emph{no} certified-by-counting neighborhood search --- around Carmichael
numbers or any other quickly-constructible markers --- can approach the
$X^{1/2}$ record, because the certificate itself forces super-linear work. A
competitive algorithm must certify success pointwise (an interval theorem such
as \cite{BHP}), or by computation (the analytic method \cite{LO}), or by the
kind of parity/algebraic certificate proposed by the Polymath project
\cite{Pol}, whose Conjecture~1.1 remains the door: deciding in time
$N^{1/2-c}$ whether an interval of length $N^{1/2+c}$ contains a prime would
break the barrier. We record one quantitative remark on the obstruction
identified there. The Polymath method computes, in time $N^{1/2-c+o(1)}$, the
parity of the number of primes in a window in every residue class modulo every
prime $r \le N^{c/10}$; it wins immediately unless every such class count is
even. The following shows that trivial counting cannot exclude that
conspiracy.

\begin{lemma}[even-pairing is cheap to sustain]
\label{lem:pairing}
Let $p_1 < \cdots < p_s$ be primes in an interval of length $\le N$, and
suppose that for every prime $r \le R$ every residue class modulo $r$
contains an even number of the $p_i$. Then $\pi(R) \le s \log_2 N$.
\end{lemma}

\begin{proof}
Fix $r$. Pairing the $p_i$ within classes gives at least $s/2$ disjoint
pairs with $r \mid p_i - p_j$. Summing over $r \le R$ counts at least
$(s/2)\pi(R)$ incidences $(\{p_i, p_j\}, r)$ with $r \mid p_i - p_j$; each of
the at most $s^2/2$ pairs contributes at most $\omega(p_i - p_j) \le \log_2 N$
incidences. Hence $(s/2)\pi(R) \le (s^2/2)\log_2 N$.
\end{proof}

A window of length $N^{1/2+c}$ has $s \le N^{1/2+c}$ primes, so
Lemma~\ref{lem:pairing} only forbids the conspiracy for
$R \gg s \log N$ --- vastly beyond the $N^{c/10}$ budget the method can
afford. Ruling it out below that range is a question about primes in
arithmetic progressions \emph{within short intervals}, where the conductor of
the associated $L$-functions consumes exactly the gain: the same wall as
Proposition~\ref{prop:barrier}, seen from a third side.

\section{The AGP pigeonhole as a prime-finding theorem}
\label{sec:agpset}

The negative results above concern primes \emph{near} Carmichael numbers. We
close the circle with a positive observation about the machinery itself: the
theorem of Alford--Granville--Pomerance that supplies the prime pool for the
companion construction (\cite{AGP}, Theorem~3.1; quoted as Lemma~2.2 of
\cite{Bro}) is, read directly, a prime-existence statement for a thin explicit
set, and hence a deterministic prime-finding algorithm. We have not found this
observation in the literature.

Fix $B < 5/12$ and let $D_B$, $z_3(B)$ be the effective constants of
\cite{AGP}, Theorem~3.1. Given $x$, set $T = \lceil D_B + 2 + \log_2 \log x
\rceil$ and let $L$ be the product of the first $T$ primes exceeding
$64T/(1-B)$. Then $L$ is squarefree, its prime factors lie below
$x^{(1-B)/2}$, $\sum_{q \mid L} 1/q \le (1-B)/32$, and $L \le x^{B}$ for
large $x$, so the theorem applies: some $k \le x^{1-B}$ with $(k, L) = 1$
satisfies
\[
\#\{\, d \mid L : dk + 1 \le x \ \text{prime} \,\}
\;\ge\; \frac{2^{-D_B-2}}{\log x}\, 2^{T} \;\ge\; 1 .
\]
Thus the explicit set
$S_x = \{\, dk+1 \le x : d \mid L,\ k \le x^{1-B},\ (k,L)=1 \,\}$, of
cardinality at most $2^{T} x^{1-B} = x^{1-B+o(1)}$, provably contains a
prime, and enumerating it with the AKS test finds one in deterministic time
$x^{1-B+o(1)}$.

As stated, the pigeonhole does not control the \emph{size} of the primes it
produces, so it finds \emph{a} prime cheaply rather than a prime exceeding a
given bound. The dyadic form needed for prime-finding proper follows,
however, directly from the actual equidistribution theorem of \cite{AGP}
(their Theorem~3.1), which is stronger than the pigeonhole consequence
quoted above: it is a \emph{two-sided} relative-error asymptotic,
\[
\Bigl| \pi(y; d, a) - \frac{\pi(y)}{\varphi(d)} \Bigr|
\;\le\; \varepsilon\, \frac{\pi(y)}{\varphi(d)},
\]
pointwise for every $d \le x^{1/A - \delta}$ not divisible by a member of a
set of at most $D_{\varepsilon,\delta}$ bad integers, valid for all
$y \ge d\, x^{1 - 1/A + \delta}$, where $A$ is any admissible zero-density
exponent (\cite{AGP}, eq.~(3.1)); \cite{AGP} note that their proof of this
statement is effective. Subtracting the asymptotic at $y = d x^{1-B}$ and at
$y = x/2$ (both in range when $B \le 1/A - \delta$ and
$d \ge 0.55\, x^{B}$) counts primes $p \equiv 1 \bmod d$ with
$p \in (x/2,\, dx^{1-B}]$, hence with cofactor $k = (p-1)/d \le x^{1-B}$;
summing over the divisors of $L$ in the dyadic block $(0.55\,x^B, x^B]$ that
avoid the bad set, and pigeonholing over $k \le x^{1-B}$, yields a shift $k$
for which some $dk + 1 \in (x/2, x]$ is prime, provided the block contains
$\gg 2^{D} \log x$ good divisors. That supply is arranged constructively:
take $T = O(D + \log\log x)$ primes in an interval $(Q, 2Q)$ with
$Q^{T/2} \approx x^{1/A - \delta}$; the $\binom{T}{T/2}$ subset products of
size $T/2$ spread over $T/2$ dyadic blocks, so some block receives
$\ge 2^{T}/T^2$ of them, and discarding divisors met by the bad set costs a
further factor $2^{-D}$ at most. (The hypotheses $(k, L) = 1$ and
$\sum_{q \mid L} 1/q$ small in the quoted pigeonhole are needed for the
Carmichael application, not here.) This proves:

\begin{theorem}[unconditional]
\label{prop:agpset}
For every fixed $B < 5/12$ there is a deterministic algorithm that, given
$N$, outputs a prime in $(N/2, N]$ in time $N^{1-B+o(1)}$; in particular in
time $N^{7/12 + o(1)}$. The constants are effectively computable, and the
exponent is $1 - 1/A + o(1)$ for any admissible zero-density exponent $A$ in
\textup{(3.1)} of \cite{AGP}.
\end{theorem}

Three remarks. First, the exponent slots into the known ladder as follows:
Lagarias--Odlyzko $N^{1/2+o(1)}$ \cite{LO}; the Baker--Harman--Pintz
interval scan $N^{0.525+o(1)}$ \cite{BHP}; Proposition~\ref{prop:agpset} at
$N^{7/12} = N^{0.5834}$; the Meissel--Lehmer binary search and the
Heath-Brown $a^3 + 2b^3$ scan, both $N^{2/3+o(1)}$ (see \cite{Pol}). The new
entry is the only one of a multiplicative species --- a union of
$\log^{O(1)} N$ values $d \mid L$ sampled along $N^{1-B}$ shifts --- and the
only one whose exponent is a dial rather than a constant. Second, the dial:
the mechanism $B < 1/A$ is stated by \cite{AGP} themselves (``if
$A \in \mathcal{A}$, then all $B < 1/A$ are in $\mathcal{B}$''), with
Huxley's $A = 12/5$ giving $5/12$. The Guth--Maynard improvement of the
zero-density exponent to $A = 30/13$ would lift $B$ to $13/30$ and the
exponent of Theorem~\ref{prop:agpset} to $N^{17/30+o(1)}$ --- their
short-interval exponent, reached by a purely multiplicative route ---
provided $30/13$ is admissible in the hybrid $q$-aspect sense of
\textup{(3.1)} of \cite{AGP}, which is proved there for the $t$-aspect and
remains to be checked for conductors. Progress toward the density
hypothesis ($A \to 2^{+}$) pushes the exponent to $N^{1/2+\varepsilon}$, at
which point an elementary scan matches the analytic record --- and there it
stops, unconditionally: \cite{AGP} observe that \textup{(3.1)} cannot hold
for any $A \le 2$, because $L(s,\chi)$ really does have $\asymp Td\log(Td)$
zeros up to height $T$ on the critical line. The infimum of the exponent
$1 - 1/A$ over \emph{all possible} zero-density inputs is therefore exactly
$1/2$: the route can match the record in the limit and can never beat it.
Nor can level-of-distribution technology rescue it. The divisor family
$\{d \mid L\}$ is necessarily sparse ($\tau(L) = N^{o(1)}$), so the
pigeonhole needs \emph{pointwise} control with bounded exceptions ---
zero-density territory, capped as above. The theorems that go beyond level
$1/2$ bound errors absolutely, at the scale of a dense family: this
includes the well-factorable theorem of Bombieri--Friedlander--%
Iwaniec, which does apply here --- indicator weights on the divisors
of a smooth $L$ are the archetype of well-factorable weights --- but its
error term $x/\log^A x$ swamps our thin main term $x^{1-B+o(1)}$; and a
relative-error average theorem for thin families at level beyond $1/2$
would, applied to a singleton family, exceed what GRH gives pointwise.
Sparse needs pointwise; pointwise caps at $1/2$; dense forfeits the gain:
the same square-root wall as Propositions~\ref{prop:barrier} and
\ref{prop:cost}, probed along yet another axis of approach and holding.
Third, the machinery is dual-use. Built to construct Carmichael numbers, it
acquires primes directly: Sections~\ref{sec:tradeoff}--\ref{sec:barriers}
close the route through proximity, and the pigeonhole at its heart is a
prime-existence theorem outright.

\section{Empirical results}
\label{sec:empirical}

The theorems above are consistent with what one actually observes.
Normalising the distance from $m$ to the nearest prime by $\log m$
(so a ``random odd number'' baseline is $\tfrac12$ in the mean):

\begin{center}
\begin{tabular}{lrrr}
\toprule
population & $n$ & mean $d/\log m$ & median \\
\midrule
all Carmichael $m \le 10^8$ (exact census) & 255 & 0.550 & 0.454 \\
\quad matched random odd composites & 12750 & 0.482 & 0.354 \\
\addlinespace
constructed Carmichael, $23$--$77$ digits & 1000 & 0.483 & 0.343 \\
\quad matched random odd & 1000 & 0.494 & 0.359 \\
\quad matched random $\equiv 1 \bmod 720720$ & 1000 & 0.508 & 0.360 \\
\bottomrule
\end{tabular}
\end{center}

\smallskip
Carmichael numbers are not closer to primes than random numbers of their size;
at small scales they are, if anything, slightly farther. What their structure
does control is \emph{where} the nearby primes sit: writing $m \equiv 1
\pmod{L}$ with $\lambda$-structure as in \cite{Bro}, a candidate $m + c$
avoids every small prime $r \mid \lambda(L)$ unless $r \mid c + 1$. In all
$1000$ constructed cases the nearest prime occupied such an admissible offset
($\gcd(c \pm 1, 3\cdot5\cdot7\cdot11\cdot13) = 1$), and a search that skips
inadmissible offsets found a certified prime above $10^{40}$ in a median of
$12$ (maximum $152$) Baillie--PSW tests per neighborhood over $1000$
engagements, without a miss. This is a constant-factor convenience, exactly as the
theory predicts: the Mertens boost $\prod_{r \mid \lambda} r/(r-1)$ on
admissible offsets is cancelled, on average, by their sparsity. Scripts
reproducing all computations in this section are available from the author.

\section{Related work and closing remarks}
\label{sec:related}

Pintz, Steiger and Szemer\'edi \cite{PSS} exhibited an infinite set of primes,
of counting density $\gg x^{2/3}/\log x$, whose membership can be decided in
deterministic polynomial time; but their \emph{generation} procedure is
randomized, and no deterministic procedure is known that produces an element
of their set (or any prime) beyond a given bound $n$ in time $n^{o(1)}$. The
companion paper \cite{Bro} achieves precisely this on the composite side, for
Carmichael numbers; the present note quantifies why that success does not
transfer to primes through proximity: the transfer is throttled by Carmichael
\emph{counting} (Theorem~\ref{thm:tradeoff} and
Corollary~\ref{cor:threshold}), capped below polynomial proximity by the
Pomerance upper bound (Proposition~\ref{prop:barrier}), and in its algorithmic
form defeated by the cost of counting certificates
(Proposition~\ref{prop:cost}).

We close by lasing the one target this analysis designates:
\emph{any improvement of Harman's $\CC(X) \ge X^{0.33}$ past $X^{0.706}$
makes Corollary~\ref{cor:threshold} nonvacuous}, yielding primes provably
closer to (infinitely many) Carmichael numbers than they are known to be to
anything else. Conversely, improvements in the gap moment \eqref{eq:sta}
lower the $0.706$ threshold toward the abundance conjecture's $\theta \to 1$;
the two literatures meet in the middle of this trade-off.

\begin{thebibliography}{99}

\bibitem{AGP} W.~R. Alford, A. Granville, C. Pomerance,
\emph{There are infinitely many Carmichael numbers},
Ann. of Math. (2) \textbf{140} (1994), 703--722.
\href{https://doi.org/10.2307/2118576}{doi:10.2307/2118576}.

\bibitem{BHP} R.~C. Baker, G. Harman, J. Pintz,
\emph{The difference between consecutive primes, II},
Proc. London Math. Soc. (3) \textbf{83} (2001), 532--562.
\href{https://doi.org/10.1112/plms/83.3.532}{doi:10.1112/plms/83.3.532}.

\bibitem{Bro} D. Brown,
\emph{A deterministic quasi-polylogarithmic algorithm for constructing
Carmichael numbers}, preprint (2026).

\bibitem{Erd} P. Erd\H{o}s,
\emph{On pseudoprimes and Carmichael numbers},
Publ. Math. Debrecen \textbf{4} (1956), 201--206.
\href{https://doi.org/10.5486/PMD.1956.4.3-4.16}{doi:10.5486/PMD.1956.4.3-4.16}.

\bibitem{GP} A. Granville, C. Pomerance,
\emph{Two contradictory conjectures concerning Carmichael numbers},
Math. Comp. \textbf{71} (2002), 883--908.
\href{https://doi.org/10.1090/S0025-5718-01-01355-2}{doi:10.1090/S0025-5718-01-01355-2}.

\bibitem{Har} G. Harman,
\emph{On the number of Carmichael numbers up to $x$},
Bull. London Math. Soc. \textbf{37} (2005), 641--650.
\href{https://doi.org/10.1112/S0024609305004686}{doi:10.1112/S0024609305004686}.

\bibitem{HB} D.~R. Heath-Brown,
\emph{The differences between consecutive primes, III},
J. London Math. Soc. (2) \textbf{20} (1979), 177--178.
\href{https://doi.org/10.1112/jlms/s2-20.2.177}{doi:10.1112/jlms/s2-20.2.177}.

\bibitem{LO} J.~C. Lagarias, A.~M. Odlyzko,
\emph{Computing $\pi(x)$: an analytic method},
J. Algorithms \textbf{8} (1987), 173--191.
\href{https://doi.org/10.1016/0196-6774(87)90037-X}{doi:10.1016/0196-6774(87)90037-X}.

\bibitem{PSS} J. Pintz, W.~L. Steiger, E. Szemer\'edi,
\emph{Infinite sets of primes with fast primality tests and quick generation
of large primes}, Math. Comp. \textbf{53} (1989), 399--406.
\href{https://doi.org/10.1090/S0025-5718-1989-0968154-X}{doi:10.1090/S0025-5718-1989-0968154-X}.

\bibitem{Pol} D.~H.~J. Polymath,
\emph{Deterministic methods to find primes},
Math. Comp. \textbf{81} (2012), 1233--1246.
\href{https://doi.org/10.1090/S0025-5718-2011-02542-1}{doi:10.1090/S0025-5718-2011-02542-1}.

\bibitem{Pom} C. Pomerance,
\emph{On the distribution of pseudoprimes},
Math. Comp. \textbf{37} (1981), 587--593.
\href{https://doi.org/10.1090/S0025-5718-1981-0628717-0}{doi:10.1090/S0025-5718-1981-0628717-0}.

\bibitem{Sel} A. Selberg,
\emph{On the normal density of primes in small intervals, and the difference
between consecutive primes}, Arch. Math. Naturvid. \textbf{47} (1943), 87--105.

\bibitem{Sta} J. Stadlmann,
\emph{On the mean square gap between primes},
\href{https://arxiv.org/abs/2212.10867}{arXiv:2212.10867} (2022).

\end{thebibliography}
\end{document}
